\documentclass[aps,prl,reprint,amsmath,amssymb,nofootinbib,floatfix]{revtex4-2}

\usepackage[utf8]{inputenc}
\usepackage[T1]{fontenc}
\usepackage{lmodern}
\usepackage{microtype}
\usepackage{amsthm}
\usepackage{mathtools}
\usepackage{booktabs}
\usepackage{xcolor}
\usepackage[
  colorlinks=true,
  linkcolor=blue!60!black,
  citecolor=blue!60!black,
  urlcolor=blue!60!black,
  pdftitle={Universal Wigner Negativity for Pure Spin States via pi-Rotation Orthogonality},
  pdfauthor={Ximing Wang, Xilin Hou}
]{hyperref}

\newtheorem{theorem}{Theorem}
\newtheorem{lemma}{Lemma}
\newtheorem{corollary}{Corollary}

\newcommand{\ket}[1]{\lvert #1 \rangle}
\newcommand{\bra}[1]{\langle #1 \rvert}
\newcommand{\braket}[2]{\langle #1 \vert #2 \rangle}
\newcommand{\Tr}{\operatorname{Tr}}
\newcommand{\divS}{\operatorname{div}_{S^2}}
\newcommand{\gradS}{\nabla_{S^2}}
\newcommand{\LapS}{\Delta_{S^2}}
\newcommand{\bn}{\mathbf{n}}
\newcommand{\bu}{\mathbf{u}}
\newcommand{\bJ}{\mathbf{J}}
\newcommand{\be}{\mathbf{e}}

\begin{document}

\title{Universal Wigner Negativity for Pure Spin States via \texorpdfstring{$\pi$}{pi}-Rotation Orthogonality}

\author{Ximing Wang}
\email{vexpaer@gmail.com}
\thanks{These authors contributed equally to this work.}
\noaffiliation

\author{Xilin Hou}
\email{1702402958@qq.com}
\thanks{These authors contributed equally to this work.}
\noaffiliation

\date{\today}

\begin{abstract}
We prove that for every spin $j \ge 1/2$ and every pure state $\ket{\psi} \in \mathbb{C}^{2j+1}$, every $\mathrm{SU}(2)$-covariant self-dual Stratonovich--Weyl Wigner function on the sphere $S^2$ strictly attains negative values. This resolves the conjecture of Davis \textit{et al.} (Open Quantum Problem F18) and completes Hudson's theorem across constant-curvature phase spaces. The proof establishes that every pure spin state is orthogonalized by a $\pi$-rotation about some axis $\bu_* \in S^2$, and is fully machine-verified in Lean~4.
\end{abstract}

\maketitle

\emph{Introduction.}---Negativity of the Wigner quasiprobability distribution is a canonical signature of quantum nonclassicality and a necessary resource for quantum computational advantage beyond classically simulable subtheories. On the flat phase space $\mathbb{R}^{2n}$ of the Heisenberg--Weyl group, Hudson's theorem~\cite{Hudson1974,SotoClaverie1983} establishes that the pure states with globally non-negative Wigner functions are precisely the Gaussian pure states. On discrete odd-dimensional phase spaces, Gross~\cite{Gross2006} proved the analogous characterization in terms of stabilizer states.

For spin-$j$ quantum systems ($j \in \frac{1}{2}\mathbb{N} := \{\frac{1}{2}, 1, \frac{3}{2}, \dots\}$), the symmetry group is $\mathrm{SU}(2)$ and the classical phase space is the sphere $S^2 \cong \mathrm{SU}(2)/\mathrm{U}(1)$ of constant positive curvature. The Stratonovich--Weyl correspondence~\cite{Stratonovich1957,VarillyGraciaBondia1989,BrifMann1999} maps each operator $A$ on $\mathcal{H}_j \cong \mathbb{C}^{2j+1}$ to a spherical Wigner function $W_A(\bn) = \Tr[A\,\Delta_j(\bn)]$ via a self-dual, $\mathrm{SU}(2)$-covariant kernel $\Delta_j(\bn)$. Studying spin coherent, cat, and Dicke states, Davis, Kumari, Mann, and Ghose~\cite{Davis2021} conjectured that \emph{no} pure spin state has vanishing Wigner negativity. Davis, Hennigar, Mann, and Ghose~\cite{Davis2023} subsequently proved strict negativity for all spin coherent states ($j \ge \frac{1}{2}$), leaving general pure states open (Open Quantum Problem F18). For mixed spin states, Denis, Davis, Mann, and Martin~\cite{Denis2024} characterized the eigenvalue polytopes of absolutely Wigner-bounded states.

Recent work on constant-curvature phase spaces brought the geometric picture into sharp focus: for the positive discrete series of $\mathrm{SU}(1,1)$ on the two-sheeted hyperboloid $H^2$ (constant negative curvature), Klimov \textit{et al.}~\cite{Klimov2026} proved that all Perelomov coherent states are strictly Wigner-positive, and Kam~\cite{Kam2026} proved that the Wigner-positive pure states strictly exceed the coherent orbit, noting that settling the $\mathrm{SU}(2)$ conjecture would complete the constant-curvature classification.

In this Letter, we prove the conjecture in full generality for all spins $j \in \frac{1}{2}\mathbb{N}$, all pure states $\ket{\psi} \in \mathcal{H}_j$, and all $2^{2j}$ self-dual $\mathrm{SU}(2)$-covariant Stratonovich--Weyl kernels. The proof centers on a geometric theorem of independent interest: every pure spin-$j$ state ($j \ge \frac{1}{2}$) has an orthogonal partner reachable by a single $\pi$-rotation in $\mathrm{SU}(2)$.

\emph{Stratonovich--Weyl framework.}---Let $D^{(j)} : \mathrm{SU}(2) \to \mathrm{U}(\mathcal{H}_j)$ be the $(2j+1)$-dimensional irreducible unitary representation with Hermitian generators $\bJ = (J_x, J_y, J_z)$ satisfying $[J_a, J_b] = i\varepsilon_{abc} J_c$ and $\bJ^2 = j(j+1)I$. For $\bn \in S^2$, let $\{\ket{j,m;\bn}\}_{m=-j}^{j}$ be an eigenbasis of $\bn\cdot\bJ$ with eigenvalues $m$, and let $g \mapsto R_g \in \mathrm{SO}(3)$ denote the double cover satisfying $D^{(j)}(g)(\bn\cdot\bJ)D^{(j)}(g)^\dagger = (R_g\bn)\cdot\bJ$.

Let $\{T^{(j)}_{lq} : 0 \le l \le 2j,\ -l \le q \le l\}$ be the orthonormal irreducible multipole operators on $\mathcal{H}_j$~\cite{VarillyGraciaBondia1989,Davis2021}, satisfying $\Tr[(T^{(j)}_{lq})^\dagger T^{(j)}_{l'q'}] = \delta_{ll'}\delta_{qq'}$, $(T^{(j)}_{lq})^\dagger = (-1)^q T^{(j)}_{l,-q}$, and $T^{(j)}_{00} = I/\sqrt{2j+1}$. For any sign vector $\boldsymbol{\epsilon} = (\epsilon_0, \epsilon_1, \dots, \epsilon_{2j})$ with $\epsilon_0 = 1$ and $\epsilon_l \in \{\pm 1\}$ ($1 \le l \le 2j$), the associated self-dual $\mathrm{SU}(2)$-covariant Stratonovich--Weyl kernel is
\begin{equation}\label{eq:SW-kernel}
\Delta_j^{(\boldsymbol{\epsilon})}(\bn)
= \sqrt{\frac{4\pi}{2j+1}} \sum_{l=0}^{2j} \epsilon_l \sum_{q=-l}^{l} Y_{lq}^*(\bn)\, T^{(j)}_{lq},
\end{equation}
where $Y_{lq}$ are the orthonormal spherical harmonics on $S^2$ with respect to the surface measure $d\Omega$. The canonical kernel~\cite{VarillyGraciaBondia1989,Davis2021,Davis2023,Kam2026} corresponds to $\epsilon_l = +1$ for all $l$.

For a pure state $\ket{\psi}$, its Wigner function is $W_\psi(\bn) := \Tr[\ket{\psi}\bra{\psi}\,\Delta_j^{(\boldsymbol{\epsilon})}(\bn)]$. With the normalized invariant measure $d\mu(\bn) := \frac{2j+1}{4\pi}\,d\Omega(\bn)$, every kernel~\eqref{eq:SW-kernel} satisfies:
(i) \emph{Reality}: $W_\psi(\bn) \in \mathbb{R}$;
(ii) \emph{Standardization}: $\int_{S^2} W_\psi(\bn)\,d\mu(\bn) = 1$;
(iii) \emph{Covariance}:
\begin{equation}\label{eq:covariance}
W_{D^{(j)}(g)\psi}(\bn) = W_\psi(R_g^{-1}\bn) \qquad \forall g \in \mathrm{SU}(2);
\end{equation}
and (iv) \emph{Traciality (self-duality)}:
\begin{equation}\label{eq:traciality}
\Tr[AB] = \int_{S^2} W_A(\bn)\,W_B(\bn)\,d\mu(\bn)
\end{equation}
for all Hermitian operators $A, B$ on $\mathcal{H}_j$. Furthermore, by Eq.~\eqref{eq:SW-kernel}, $W_\psi$ is a finite linear combination of spherical harmonics of degrees $0 \le l \le 2j$, hence the restriction to $S^2$ of a real polynomial of degree at most $2j$ on $\mathbb{R}^3$.

\emph{Orthogonality obstruction.}---We first show that having any orthogonal partner in the $\mathrm{SU}(2)$ orbit forbids global non-negativity of $W_\psi$.

\begin{lemma}[Orthogonal orbit partner]\label{lem:ortho-neg}
Let $j \in \frac{1}{2}\mathbb{N}$ and let $\ket{\psi} \in \mathcal{H}_j$ be a unit vector. If there exists $g \in \mathrm{SU}(2)$ such that $\bra{\psi} D^{(j)}(g) \ket{\psi} = 0$, then for every self-dual kernel~\eqref{eq:SW-kernel}, $\min_{\bn \in S^2} W_\psi(\bn) < 0$.
\end{lemma}

\begin{proof}
Suppose for contradiction that $W_\psi(\bn) \ge 0$ for all $\bn \in S^2$. Setting $A = \ket{\psi}\bra{\psi}$ and $B = D^{(j)}(g)\ket{\psi}\bra{\psi}D^{(j)}(g)^\dagger$ in Eqs.~\eqref{eq:covariance}--\eqref{eq:traciality} gives
\begin{equation}\label{eq:zero-overlap}
0 = \big|\bra{\psi} D^{(j)}(g) \ket{\psi}\big|^2 = \int_{S^2} W_\psi(\bn)\,W_\psi(R_g^{-1}\bn)\,d\mu(\bn).
\end{equation}
By continuity and non-negativity, $W_\psi(\bn)\,W_\psi(R_g^{-1}\bn) = 0$ for all $\bn \in S^2$. Since $\int_{S^2} W_\psi\,d\mu = 1 > 0$, there is a nonempty open set $U \subset S^2$ on which $W_\psi > 0$; thus $W_\psi \equiv 0$ on the nonempty open set $V := R_g^{-1}(U)$.

Fix $\mathbf{p} \in V$. Any $\mathbf{q} \in S^2$ lies on a great circle $\gamma(t) = (\cos t)\mathbf{p} + (\sin t)\mathbf{w}$ ($\mathbf{w} \perp \mathbf{p}$, $|\mathbf{w}|=1$). Since $W_\psi$ is a polynomial of degree at most $2j$ on $\mathbb{R}^3$, $f(t) := W_\psi(\gamma(t)) = \sum_{k=-2j}^{2j} c_k e^{ikt}$ is a trigonometric polynomial vanishing on an open interval $(-\varepsilon, \varepsilon)$ around $0$. The algebraic polynomial $P(z) := z^{2j}\sum_{k=-2j}^{2j} c_k z^k$ of degree at most $4j$ therefore has infinitely many roots on the unit circle, forcing $P \equiv 0$ and $W_\psi(\mathbf{q}) = 0$. Hence $W_\psi \equiv 0$, contradicting $\int_{S^2} W_\psi\,d\mu = 1$.
\end{proof}

\emph{Two structural lemmas for integer spins.}---To locate an orthogonalizing $\pi$-rotation for integer spins $j \in \mathbb{N}$, we combine two short results.

\begin{lemma}[Zeros of complex spherical harmonics]\label{lem:complex-sh-zero}
Let $Y \in C^\infty(S^2, \mathbb{C})$ be a nonzero complex eigenfunction of the Laplace--Beltrami operator $\LapS$ with eigenvalue $-\lambda < 0$: $\LapS Y = -\lambda Y$. Then there exists $\bn_* \in S^2$ such that $Y(\bn_*) = 0$.
\end{lemma}

\begin{proof}
Suppose $Y(\bn) \neq 0$ for all $\bn \in S^2$. Write $Y = u + iv$ and $r := |Y| = \sqrt{u^2 + v^2} > 0$. Because $S^2$ is simply connected, the smooth map $Y/r : S^2 \to S^1$ admits a smooth global phase lift $\theta \in C^\infty(S^2, \mathbb{R})$ with $u = r\cos\theta$ and $v = r\sin\theta$. The tangent vector field $\mathbf{V} := u\gradS v - v\gradS u$ satisfies $\mathbf{V} = r^2\gradS\theta$, while $\LapS u = -\lambda u$ and $\LapS v = -\lambda v$ yield
\begin{equation}
\divS(\mathbf{V}) = u\LapS v - v\LapS u = u(-\lambda v) - v(-\lambda u) = 0.
\end{equation}
Applying the divergence theorem on closed $S^2$ to $\theta r^2\gradS\theta$ gives
\begin{equation}
0 = \int_{S^2} \divS(\theta r^2\gradS\theta)\,d\Omega = \int_{S^2} r^2 |\gradS\theta|^2\,d\Omega.
\end{equation}
Since $r > 0$ and $S^2$ is connected, $\theta \equiv \theta_0$ is constant, so $e^{-i\theta_0}Y = r > 0$ everywhere and $\int_{S^2} e^{-i\theta_0}Y\,d\Omega > 0$. Yet $\lambda > 0$ implies $\int_{S^2} Y\,d\Omega = -\lambda^{-1}\int_{S^2}\LapS Y\,d\Omega = 0$, a contradiction.
\end{proof}

Note that Lemma~\ref{lem:complex-sh-zero} holds identically on any closed, connected, simply connected Riemannian manifold, and simple connectivity is sharp (on $S^1$, $Y(t) = e^{it}$ has $\lambda = 1$ and no zeros).

\begin{table*}[t]
\caption{\label{tab:trichotomy}Classification of pure-state Wigner positivity across two-dimensional homogeneous phase spaces of constant Gaussian curvature $\kappa$.}
\begin{ruledtabular}
\begin{tabular}{llll}
Curvature $\kappa$ & Symmetry Group & Phase Space & Wigner-Positive Pure States \\
\colrule
$\kappa < 0$ & $\mathrm{SU}(1,1)$ & Hyperboloid $H^2$ & Strictly larger than Perelomov coherent orbit~\cite{Klimov2026,Kam2026} \\
$\kappa = 0$ & Heisenberg--Weyl & Plane $\mathbb{R}^2$ & Exactly Gaussian / coherent orbit~\cite{Hudson1974} \\
$\kappa > 0$ & $\mathrm{SU}(2)$ & Sphere $S^2$ & Empty ($\emptyset$) for all $j \in \frac{1}{2}\mathbb{N}$ (Theorem~\ref{thm:main}) \\
\end{tabular}
\end{ruledtabular}
\end{table*}

\begin{lemma}[Equatorial integral identity]\label{lem:equatorial}
Let $j \in \mathbb{N}$. For any $\bn \in S^2$, pick a right-handed orthonormal frame $(\be_1, \be_2, \bn)$ and set $\bu(\phi) := \cos\phi\,\be_1 + \sin\phi\,\be_2$. Then
\begin{equation}\label{eq:equatorial}
\frac{1}{2\pi}\int_{0}^{2\pi} e^{-i\pi \bu(\phi)\cdot\bJ}\,d\phi = (-1)^j \ket{j,0;\bn}\bra{j,0;\bn}.
\end{equation}
\end{lemma}

\begin{proof}
By $\mathrm{SO}(3)$-covariance of $D^{(j)}$ for integer $j$, it suffices to take $(\be_1, \be_2, \bn) = (\hat{\mathbf{x}}, \hat{\mathbf{y}}, \hat{\mathbf{z}})$. Since $e^{-i\phi J_z} J_x e^{i\phi J_z} = \cos\phi\,J_x + \sin\phi\,J_y$, we have $e^{-i\pi\bu(\phi)\cdot\bJ} = e^{-i\phi J_z} e^{-i\pi J_x} e^{i\phi J_z}$. Integrating matrix elements in the $J_z$-basis gives
\begin{align}
\frac{1}{2\pi}\int_{0}^{2\pi} &\bra{j,m'} e^{-i\phi J_z} e^{-i\pi J_x} e^{i\phi J_z} \ket{j,m}\,d\phi \nonumber\\
&= \delta_{m',m}\,\bra{j,m} e^{-i\pi J_x} \ket{j,m}.
\end{align}
Because $e^{-i\pi J_x} J_z e^{i\pi J_x} = -J_z$, the operator $e^{-i\pi J_x}$ maps the $m$-eigenspace to the $(-m)$-eigenspace, so every diagonal element with $m \neq 0$ vanishes. For $m = 0$, writing $e^{-i\pi J_x} = e^{-i\frac{\pi}{2}J_z} e^{-i\pi J_y} e^{i\frac{\pi}{2}J_z}$ and using $J_z\ket{j,0}=0$ yields $\bra{j,0}e^{-i\pi J_x}\ket{j,0} = d^{(j)}_{0,0}(\pi) = P_j(-1) = (-1)^j$, proving Eq.~\eqref{eq:equatorial}.
\end{proof}

\emph{Main theorems.}---Combining Lemmas~\ref{lem:complex-sh-zero} and~\ref{lem:equatorial} establishes universal $\pi$-rotation orthogonality.

\begin{theorem}[$\pi$-rotation orthogonality]\label{thm:pi-ortho}
For every spin $j \in \frac{1}{2}\mathbb{N}$ ($j \ge \frac{1}{2}$) and every pure state $\ket{\psi} \in \mathcal{H}_j$, there exists a unit axis $\bu_* \in S^2$ such that
\begin{equation}\label{eq:pi-ortho}
\bra{\psi} e^{-i\pi \bu_* \cdot \bJ} \ket{\psi} = 0.
\end{equation}
\end{theorem}

\begin{proof}
\emph{Case 1: Half-integer spin $j \in \mathbb{N}_0 + \frac{1}{2}$.} For any $\bu \in S^2$, $\bu\cdot\bJ$ has eigenvalues $m = k + \frac{1}{2}$ ($k \in \mathbb{Z}$), so $H(\bu) := i\,e^{-i\pi \bu\cdot\bJ}$ has real eigenvalues $i\,e^{-i\pi(k+1/2)} = (-1)^k \in \{\pm 1\}$ and is Hermitian. Moreover, sending $\bu \mapsto -\bu$ reverses $m \mapsto -m$ and flips $i\,e^{-i\pi(-m)} = -i\,e^{-i\pi m}$, so $H(-\bu) = -H(\bu)$. Thus $h(\bu) := \bra{\psi} H(\bu) \ket{\psi}$ is continuous, real-valued, and odd on connected $S^2$. By the intermediate value theorem along any path from $\bu_0$ to $-\bu_0$, there exists $\bu_* \in S^2$ with $h(\bu_*) = 0$.

\emph{Case 2: Integer spin $j \in \mathbb{N}$.} Since $j \in \mathbb{Z}$, $D^{(j)}$ descends to $\mathrm{SO}(3)$ and $e^{-i\alpha J_z}\ket{j,0} = \ket{j,0}$, so $\ket{j,0;\bn} := D^{(j)}(R)\ket{j,0}$ (with $R\hat{\mathbf{z}} = \bn$) is uniquely defined with no phase ambiguity. Define $Y_\psi(\bn) := \braket{j,0;\bn}{\psi}$. The linear map $\Phi(\psi) := Y_\psi$ intertwines the irreducible action $D^{(j)}$ with the left regular action $(L_R f)(\bn) = f(R^{-1}\bn)$ on $C^\infty(S^2)$ and satisfies $\Phi(\ket{j,0})(\hat{\mathbf{z}}) = 1 \neq 0$; by irreducibility, $\ker(\Phi) = \{0\}$, so $Y_\psi \not\equiv 0$. Since $\Phi$ intertwines the Casimir $\bJ^2 = j(j+1)I$ with $-\LapS$, we have $\LapS Y_\psi = -j(j+1)Y_\psi$ with $\lambda = j(j+1) > 0$. By Lemma~\ref{lem:complex-sh-zero}, there exists $\bn_* \in S^2$ with $\braket{j,0;\bn_*}{\psi} = 0$.

In an orthonormal right-handed frame $(\be_1, \be_2, \bn_*)$, let $\bu(\phi) = \cos\phi\,\be_1 + \sin\phi\,\be_2$. Because $j \in \mathbb{Z}$, each $e^{-i\pi\bu(\phi)\cdot\bJ}$ has eigenvalues $(-1)^m \in \{\pm 1\}$ and is Hermitian; thus $g(\phi) := \bra{\psi}e^{-i\pi\bu(\phi)\cdot\bJ}\ket{\psi}$ is continuous and real-valued on $[0, 2\pi]$. Taking the expectation value of Lemma~\ref{lem:equatorial} at $\bn = \bn_*$ gives
\begin{equation}
\frac{1}{2\pi}\int_{0}^{2\pi} g(\phi)\,d\phi = (-1)^j \big|\braket{j,0;\bn_*}{\psi}\big|^2 = 0.
\end{equation}
Hence $g(\phi_*) = 0$ for some $\phi_* \in [0, 2\pi]$, and $\bu_* := \bu(\phi_*)$ satisfies Eq.~\eqref{eq:pi-ortho}.
\end{proof}

Notice that for half-integer spins and the canonical kernel ($\epsilon_l = +1$), Kramers' time-reversal orthogonality $\braket{\psi}{\Theta\psi} = 0$ also implies Wigner negativity via $W_{\Theta\psi}(\bn) = W_\psi(-\bn)$; using $\pi$-rotations avoids antiunitary maps and treats all $2^{2j}$ kernels~\eqref{eq:SW-kernel} uniformly.

\begin{theorem}[Universal Wigner negativity]\label{thm:main}
For every spin $j \in \frac{1}{2}\mathbb{N}$, every pure state $\ket{\psi} \in \mathcal{H}_j$, and every self-dual $\mathrm{SU}(2)$-covariant Stratonovich--Weyl kernel~\eqref{eq:SW-kernel},
\begin{equation}
\min_{\bn \in S^2} W_\psi(\bn) < 0,
\end{equation}
and the integrated Wigner negativity $\delta(\psi) := \frac{1}{2}\big(\int_{S^2} |W_\psi|\,d\mu - 1\big)$ satisfies $\delta(\psi) > 0$.
\end{theorem}

\begin{proof}
By Theorem~\ref{thm:pi-ortho}, $g_* := \exp(-i\frac{\pi}{2}\bu_*\cdot\boldsymbol{\sigma}) \in \mathrm{SU}(2)$ satisfies $\bra{\psi}D^{(j)}(g_*)\ket{\psi} = 0$. Lemma~\ref{lem:ortho-neg} then gives $\min_{S^2} W_\psi < 0$. Continuity yields $W_\psi < 0$ on a nonempty open set, so $\int_{S^2}|W_\psi|\,d\mu > \int_{S^2} W_\psi\,d\mu = 1$ and $\delta(\psi) > 0$.
\end{proof}

\begin{corollary}[Uniform fixed-spin lower bound]\label{cor:uniform}
For every fixed spin $j \in \frac{1}{2}\mathbb{N}$ and every self-dual kernel~\eqref{eq:SW-kernel},
\begin{equation}
\sup_{\ket{\psi}\in\mathbb{C}\mathrm{P}^{2j}} \min_{\bn\in S^2} W_\psi(\bn) < 0,
\qquad
\inf_{\ket{\psi}\in\mathbb{C}\mathrm{P}^{2j}} \delta(\psi) > 0.
\end{equation}
Moreover, writing $W_\psi = W_+ - W_-$ with $W_\pm := \max(\pm W_\psi, 0)$, letting $R_* := R_{\bu_*}(\pi)$, and setting $M_j := \|\Delta_j^{(\boldsymbol{\epsilon})}(\bn)\|_{\mathrm{op}}$,
\begin{equation}\label{eq:overlap-bound}
\delta(\psi) \ge \frac{1}{2M_j} \int_{S^2} W_+(\bn)\,W_+(R_*\bn)\,d\mu(\bn).
\end{equation}
\end{corollary}

\begin{proof}
Since $|W_\psi(\bn) - W_\phi(\bn)| \le M_j\|\ket{\psi}\bra{\psi} - \ket{\phi}\bra{\phi}\|_1$ uniformly on $S^2$, both $\min_{S^2} W_\psi$ and $\delta(\psi)$ are continuous on compact $\mathbb{C}\mathrm{P}^{2j}$ and attain their strictly negative/positive extrema by Theorem~\ref{thm:main}. For Eq.~\eqref{eq:overlap-bound}, expanding $0 = \int_{S^2}(W_+ - W_-)(W_+\circ R_* - W_-\circ R_*)\,d\mu$ with $R_*^2 = I$ and $R_*$-invariance of $d\mu$ gives
\begin{align}
2\!\int_{S^2}\! W_+ (W_-\!\circ R_*)\,d\mu
&= \!\int_{S^2}\! W_+(W_+\!\circ R_*)\,d\mu \nonumber\\
&\quad + \!\int_{S^2}\! W_-(W_-\!\circ R_*)\,d\mu \nonumber\\
&\ge \!\int_{S^2}\! W_+(W_+\!\circ R_*)\,d\mu.
\end{align}
Bounding $W_+ \le M_j$ on the left-hand side and using $\int_{S^2} W_-\circ R_*\,d\mu = \delta(\psi)$ yields Eq.~\eqref{eq:overlap-bound}.
\end{proof}

\emph{Discussion.}---Theorem~\ref{thm:main} completes the classification of pure-state Wigner positivity across all two-dimensional constant-curvature phase spaces (Table~\ref{tab:trichotomy}). The contrast with non-positive curvature is geometric: on $\mathbb{R}^2$, displaced coherent states never become orthogonal ($\braket{0}{\alpha} = e^{-|\alpha|^2/2} \neq 0$), and likewise no $\mathrm{SU}(1,1)$ element orthogonalizes the lowest-weight coherent state on $H^2$; on compact $S^2$, Theorem~\ref{thm:pi-ortho} shows that every pure state admits an orthogonal partner via a single $\pi$-rotation.

Two quantitative questions remain open: (i) proving the conjecture of Davis \textit{et al.}~\cite{Davis2021,Davis2023} that spin coherent states minimize $\delta(\psi)$ among all pure spin-$j$ states, and (ii) determining the minimum rotation angle $\theta_*(\psi) \le \pi$ required to orthogonalize a general pure spin-$j$ state.

To guarantee complete rigor, the entire proof chain (Lemmas~\ref{lem:ortho-neg}--\ref{lem:equatorial}, Theorems~\ref{thm:pi-ortho}--\ref{thm:main}, Corollary~\ref{cor:uniform}, and an explicit spin-$\frac{1}{2}$ witness certifying non-vacuity) has been machine-verified in Lean~4 (\texttt{v4.36.0} + \texttt{Mathlib4}, zero \texttt{sorry}, standard Lean axioms only) and cross-checked via exact-Racah numerical calculations; full verification details appear in the End Matter.

\emph{Data and Code Availability.}---The Lean~4 formalization, Python verification script, and manuscript sources are openly archived on Zenodo~\cite{Wang2026Zenodo} and GitHub at \url{https://github.com/vexpaer/su2-wigner-negativity}.

\emph{Acknowledgments and AI Use Disclosure.}---In accordance with APS editorial policies on the appropriate use of AI tools, the authors disclose that Google DeepMind Antigravity was used under the authors' direction during proof exploration, Lean~4 formalization, numerical verification scripting, and manuscript drafting. All mathematical statements and derivations were independently machine-verified without gaps in Lean~4 (\texttt{v4.36.0} with \texttt{Mathlib4}) and cross-checked via exact rational Clebsch--Gordan computations, and the authors assume full responsibility for the accuracy and integrity of the work.

\begin{thebibliography}{12}

\bibitem{Hudson1974}
R.~L.~Hudson, When is the Wigner quasi-probability density non-negative?, Rep.\ Math.\ Phys.\ \textbf{6}, 249 (1974).

\bibitem{SotoClaverie1983}
F.~Soto and P.~Claverie, When is the Wigner function of multidimensional systems nonnegative?, J.\ Math.\ Phys.\ \textbf{24}, 97 (1983).

\bibitem{Gross2006}
D.~Gross, Hudson's theorem for finite-dimensional quantum systems, J.\ Math.\ Phys.\ \textbf{47}, 122107 (2006).

\bibitem{Stratonovich1957}
R.~L.~Stratonovich, On distributions in representation space, Sov.\ Phys.\ JETP \textbf{4}, 891 (1957).

\bibitem{VarillyGraciaBondia1989}
J.~C.~V\'arilly and J.~M.~Gracia-Bond\'{\i}a, The Moyal representation for spin, Ann.\ Phys.\ (N.Y.) \textbf{190}, 107 (1989).

\bibitem{BrifMann1999}
C.~Brif and A.~Mann, Phase-space formulation of quantum mechanics and quantum-state reconstruction for physical systems with Lie-group symmetries, Phys.\ Rev.\ A \textbf{59}, 971 (1999).

\bibitem{Davis2021}
J.~Davis, M.~Kumari, R.~B.~Mann, and S.~Ghose, Wigner negativity in spin-$j$ systems, Phys.\ Rev.\ Research \textbf{3}, 033134 (2021).

\bibitem{Davis2023}
J.~Davis, R.~Hennigar, R.~B.~Mann, and S.~Ghose, Stellar representation of extremal Wigner-negative spin states, J.\ Phys.\ A: Math.\ Theor.\ \textbf{56}, 265302 (2023).

\bibitem{Denis2024}
J.~Denis, J.~Davis, R.~B.~Mann, and J.~Martin, Polytopes of Absolutely Wigner Bounded Spin States, Quantum \textbf{8}, 1550 (2024).

\bibitem{Klimov2026}
A.~B.~Klimov, A.~Mu\~noz, G.~Leuchs, J.-P.~Gazeau, and L.~L.~S\'anchez-Soto, Wigner negativity and stellar rank for $\mathrm{SU}(1,1)$ states, arXiv:2607.22810 (2026).

\bibitem{Kam2026}
C.-F.~Kam, Hudson's theorem fails for the $\mathrm{SU}(1,1)$ discrete series, arXiv:2609.07242 (2026).

\bibitem{Wang2026Zenodo}
X.~Wang, Universal Wigner Negativity for Pure Spin States via $\pi$-Rotation Orthogonality (code and Lean~4 formalization, v1.0.0), Zenodo, \url{https://doi.org/10.5281/zenodo.24563264} (2026).

\end{thebibliography}

\onecolumngrid
\vspace{12pt}
\begin{center}
\textbf{\large End Matter}
\end{center}
\vspace{4pt}
\twocolumngrid

\emph{Appendix A: Machine-Checked Formalization in Lean~4.}---The entire mathematical argument of this Letter has been formalized in Lean~4 ({\footnotesize\texttt{v4.36.0}} with {\footnotesize\texttt{Mathlib4}}), totaling $4{,}083$ lines across four self-contained modules archived in~\cite{Wang2026Zenodo}. Table~\ref{tab:lean-map} summarizes the correspondence between paper statements and verified Lean declarations.

\begin{table}[h]
\caption{\label{tab:lean-map}Correspondence between manuscript results and machine-verified Lean~4 declarations in namespace {\footnotesize\texttt{SU2Wigner}}.}
\footnotesize
\begin{ruledtabular}
\begin{tabular}{ll}
Manuscript Result & Verified Lean~4 Declaration \\
\colrule
Spin-$N/2$ rep.\ $D^{(j)}$ & \texttt{spinRep}, \texttt{spinRep\_mul} \\
Lemma~\ref{lem:ortho-neg} & \texttt{wigner\_neg\_of\_ortho\_partner} \\
Lemma~\ref{lem:complex-sh-zero} & \texttt{spherical\_harmonic\_has\_zero} \\
Lemma~\ref{lem:equatorial} & \texttt{equatorial\_integral\_identity} \\
Theorem~\ref{thm:pi-ortho} ($j\in\mathbb{N}_0+\frac{1}{2}$) & \texttt{pi\_ortho\_odd} \\
Theorem~\ref{thm:pi-ortho} ($j\in\mathbb{N}$) & \texttt{pi\_ortho\_even} \\
Theorem~\ref{thm:pi-ortho} (all $j\ge\frac{1}{2}$) & \texttt{pi\_ortho\_all} \\
Theorem~\ref{thm:main} & \texttt{universal\_wigner\_negativity} \\
Corollary~\ref{cor:uniform} & \texttt{uniform\_wigner\_negativity\_bound} \\
Non-vacuity ($j=\frac{1}{2}$) & \texttt{spinHalfStratonovichWeyl} \\
\end{tabular}
\end{ruledtabular}
\end{table}

Briefly, {\footnotesize\texttt{SpinRep.lean}} constructs the explicit unitary $(N+1)$-dimensional representation {\footnotesize\texttt{spinRep N}} of $\mathrm{SU}(2)$ via symmetric homogeneous polynomials and proves the group homomorphism law, unitarity, and {\footnotesize\texttt{pi\_ortho\_odd}}. {\footnotesize\texttt{SphericalHarmonicZero.lean}} proves Lemma~\ref{lem:complex-sh-zero} for all $C^2$ complex Laplace eigenfunctions on $S^2$ with $\lambda > 0$ using 1D calculus, Schwarz mixed-partial symmetry, meridian phase lifting, and Fubini's theorem in spherical coordinates. {\footnotesize\texttt{EquatorialIntegral.lean}} proves that $Y_\psi(\bn) = \braket{j,0;\bn}{\psi}$ is a nonvanishing spherical harmonic of eigenvalue $j(j+1)$, establishes Lemma~\ref{lem:equatorial}, and proves {\footnotesize\texttt{pi\_ortho\_even}} and {\footnotesize\texttt{pi\_ortho\_all}} for all $N = 2j \ge 1$. Finally, {\footnotesize\texttt{WignerNegativity.lean}} formalizes trigonometric-polynomial root counting on great circles, defines the axiomatic structure {\footnotesize\texttt{StratonovichWeylSystem N}}, proves Lemma~\ref{lem:ortho-neg}, Theorem~\ref{thm:main}, and Corollary~\ref{cor:uniform}, and constructs the concrete spin-$\frac{1}{2}$ kernel {\footnotesize\texttt{spinHalfStratonovichWeyl}} verifying all four axioms so that the axiomatic hypotheses are certified non-vacuous.

The full Lean build ({\footnotesize\texttt{lake build}}) completes with zero {\footnotesize\texttt{sorry}} placeholders and zero custom axioms; running {\footnotesize\texttt{\#print axioms}} on every declaration in Table~\ref{tab:lean-map} confirms exclusive dependence on Lean's three standard logical axioms ({\footnotesize\texttt{propext}}, {\footnotesize\texttt{Classical.choice}}, {\footnotesize\texttt{Quot.sound}}).

\emph{Appendix B: Independent Numerical Verification.}---As an independent sanity check on all matrix conventions and normalization constants, the archived script {\footnotesize\texttt{verify\_theorem.py}}~\cite{Wang2026Zenodo} constructs the multipole operators $T^{(j)}_{lq}$ using exact rational Racah Clebsch--Gordan formulas and directly verifies: (i) reality, standardization, covariance, and traciality of $\Delta_j^{(\boldsymbol{\epsilon})}(\bn)$ for both canonical ($\epsilon_l = +1$) and non-canonical sign vectors $\boldsymbol{\epsilon}$; (ii) the equatorial integral identity~\eqref{eq:equatorial} to machine precision; (iii) explicit zeros $\bu_* \in S^2$ of $\bra{\psi}e^{-i\pi\bu\cdot\bJ}\ket{\psi}$; and (iv) strict negativity $\min_{S^2} W_\psi < 0$ across spin coherent, Dicke, cat, and Haar-random pure states for all spins $j \in \{\frac{1}{2}, 1, \frac{3}{2}, \dots, 4\}$.

\end{document}
